\documentclass[11pt]{article}

\usepackage[margin=1in]{geometry}
\usepackage{graphicx}
\usepackage{array}
\usepackage{hyperref}
\usepackage{enumitem}
\usepackage{booktabs}
\usepackage{caption}
\usepackage{amsmath}
\usepackage{titlesec}
\usepackage{float}

\titleformat{\section}{\large\bfseries}{\thesection.}{1em}{}

\begin{document}

\begin{center}
    \Large \textbf{Sri Sivasubramaniya Nadar College of Engineering, Chennai} \\
    \large (An Autonomous Institution affiliated to Anna University) \\
    \vspace{0.3cm}
\end{center}

\begin{table}[h!]
\renewcommand{\arraystretch}{1.5}
\centering
\resizebox{\textwidth}{!}{%
\begin{tabular}{|l|cll|}
\hline
\textbf{Degree \& Branch} & \multicolumn{1}{c|}{B.E Computer Science \& Engineering} & \textbf{Semester} & VI \\ \hline
\textbf{Subject Code \& Name} & \multicolumn{3}{c|}{UCS2612 -- Machine Learning Laboratory} \\ \hline
\textbf{Academic Year} & \multicolumn{1}{c|}{2025--2026 (Even)} & \textbf{Batch} & 2023--2027 \\ \hline
\textbf{Name} & \multicolumn{1}{c|}{Mehanth T} & \textbf{Register No.} & 3122235001080 \\ \hline
\end{tabular}
}
\end{table}

\begin{center}
    \textbf{\Large Experiment 6}
\end{center}

\begin{center}
    \textbf{Decision Tree and Random Forest: A Comparative Classification Study}
\end{center}

\section*{Objective}
\begin{itemize}
    \item To implement a \textbf{Decision Tree classifier}.
    \item To extend the Decision Tree into a \textbf{Random Forest ensemble model}.
    \item To study the impact of hyperparameters on overfitting and generalization.
    \item To \textbf{select optimal hyperparameters using 5-Fold Cross-Validation}.
    \item To compare single-tree and ensemble-tree models.
\end{itemize}

\section*{Dataset}
\textbf{Wisconsin Diagnostic Breast Cancer Dataset} (loaded from \texttt{sklearn.datasets})
\begin{itemize}
    \item Total samples: 569 (357 Benign, 212 Malignant)
    \item Features: 30 numerical attributes
    \item Target classes: Malignant (0) and Benign (1)
    \item Train/Test split: 80\%/20\% (455 train, 114 test)
\end{itemize}

\section*{Theory}

\subsection*{Decision Tree Classifier}
A Decision Tree is a supervised learning model that recursively splits the feature space using impurity measures (Gini or Entropy) to form decision rules. It is prone to high variance and overfitting at large depths.

\subsection*{Random Forest Classifier}
Random Forest is an ensemble of decision trees trained on bootstrapped samples with random feature selection. Majority voting aggregates predictions, reducing variance and improving generalization.

\section*{Steps for Implementation}
\begin{enumerate}[label=\arabic*.]
    \item Load dataset and verify class distribution.
    \item Perform EDA: class distribution plot and correlation heatmap.
    \item 80--20 train/test split with stratification.
    \item Decision Tree: 5-Fold CV over criterion $\times$ max\_depth $\times$ min\_samples\_split.
    \item Select best DT hyperparameters and retrain on full train set.
    \item Random Forest: 5-Fold CV over n\_estimators $\times$ max\_depth $\times$ max\_features.
    \item Select best RF hyperparameters and retrain.
    \item Evaluate both on the held-out test set.
\end{enumerate}

\section*{Hyperparameter Tuning Results}

\subsection*{Decision Tree Cross-Fold Results}
\begin{table}[H]
\centering
\caption{Decision Tree -- 5-Fold CV Results}
\begin{tabular}{cccc}
\toprule
Criterion & Max Depth & Avg CV Accuracy (\%) & Avg CV F1 Score \\
\midrule
Gini     & 3       & 92.53 & 0.9400 \\
Gini     & 5       & 91.87 & 0.9335 \\
Entropy  & 5       & \textbf{92.97} & \textbf{0.9439} \\
Entropy  & None    & 92.97 & 0.9440 \\
Gini     & 7       & 91.87 & 0.9333 \\
\bottomrule
\end{tabular}
\end{table}

\textbf{Best DT:} criterion=\texttt{entropy}, max\_depth=5, min\_samples\_split=2

\subsection*{Random Forest Cross-Fold Results}
\begin{table}[H]
\centering
\caption{Random Forest -- 5-Fold CV Results}
\begin{tabular}{ccccc}
\toprule
n\_estimators & Max Depth & Max Features & Avg CV Accuracy (\%) & Avg CV F1 Score \\
\midrule
50  & 5    & sqrt & 95.60 & 0.9647 \\
100 & 10   & sqrt & 96.26 & 0.9699 \\
100 & None & sqrt & 96.26 & 0.9699 \\
200 & None & sqrt & 96.26 & 0.9700 \\
100 & None & log2 & \textbf{96.48} & \textbf{0.9718} \\
\bottomrule
\end{tabular}
\end{table}

\textbf{Best RF:} n\_estimators=100, max\_depth=None, max\_features=\texttt{log2}

\section*{5-Fold Cross-Validation Performance Comparison}

\begin{table}[H]
\centering
\caption{5-Fold CV Accuracy per Fold (Best Models)}
\begin{tabular}{lcccccc}
\toprule
Model & Fold 1 & Fold 2 & Fold 3 & Fold 4 & Fold 5 & Average \\
\midrule
Decision Tree & 92.98\% & 91.23\% & 97.37\% & 94.74\% & 92.92\% & 93.85\% \\
Random Forest & 96.49\% & 95.61\% & 95.61\% & 95.61\% & 96.46\% & 95.96\% \\
\bottomrule
\end{tabular}
\end{table}

\section*{Final Test-Set Evaluation}

\begin{table}[H]
\centering
\caption{Test-Set Performance (114 samples)}
\begin{tabular}{lcccccc}
\toprule
\textbf{Model} & \textbf{Accuracy} & \textbf{Precision} & \textbf{Recall} & \textbf{F1-Score} & \textbf{AUC-ROC} \\
\midrule
Decision Tree  & 92.98\% & 0.9706 & 0.9167 & 0.9429 & 0.9435 \\
Random Forest  & \textbf{95.61\%} & \textbf{0.9589} & \textbf{0.9722} & \textbf{0.9655} & \textbf{0.9924} \\
\bottomrule
\end{tabular}
\end{table}

\section*{Observation Questions}
\begin{itemize}
    \item \textbf{How does tree depth affect overfitting?} Shallow trees (depth=3--5) generalise better. Unconstrained depth (=None) fits training data perfectly but may overfit; here, with only 30 features, the effect is mild because the data is not noisy.
    \item \textbf{Which hyperparameter had the greatest impact?} The choice of \texttt{max\_features} in Random Forest (\texttt{log2} vs. \texttt{sqrt}) marginally improved accuracy. For Decision Tree, \texttt{criterion=entropy} with \texttt{max\_depth=5} gave the best balance.
    \item \textbf{How does Random Forest improve generalisation?} By training 100 trees on different bootstrap samples with random feature subsets, variance is reduced via averaging, leading to a 2.6\% accuracy gain over a single tree.
    \item \textbf{Did ensemble learning always improve performance?} Yes -- Random Forest consistently outperformed the single Decision Tree across all 5 folds and on the final test set.
\end{itemize}

\section*{Conclusion}
Decision Tree (entropy, depth=5) achieved 92.98\% test accuracy; Random Forest (100 trees, log2 features) achieved 95.61\% with an AUC of 0.9924 -- a substantial improvement demonstrating the power of ensemble averaging for variance reduction.

\section*{References}
\begin{itemize}
    \item \href{https://scikit-learn.org/stable/modules/tree.html}{Scikit-learn: Decision Trees}
    \item \href{https://scikit-learn.org/stable/modules/ensemble.html}{Scikit-learn: Random Forest}
    \item \href{https://archive.ics.uci.edu/dataset/17/breast+cancer+wisconsin+diagnostic}{UCI Dataset}
\end{itemize}

\end{document}
